\section{Conclusion}\label{sec:conclusion}
The complexity of a scalar Newton quantity need not equal that of the
Newton direction. Classical local estimation of a positive inverse form
to relative error has a variance factor linear in the condition number,
while separate sparse families with full $\SQ$ access show that the
statistical factor and the exponential factor are each necessary.
Explicit composition shows when both factors apply to the same scalar, and why the baseline of a
hidden signal cannot be ignored. Under plain block access, the known
variable-time algorithm is optimal up to logarithmic factors, but the
lower bound does not hold when exact sparse entries can be queried.

In our optimization constructions, hardness can sit in an ordinary optimal value, an optimizer coordinate,
or a small Newton decrement, with all access to the formulation
simulated. It can also add up over a sequence of requested scalars or in
one endpoint coordinate. The same examples show the limits: the optimal
value without the objective tilt may be public, a normalized state may
lose a hard sign or scale, and one hard Newton direction may be reused
along a long trajectory. Comparisons of whole optimization methods must
therefore account for input acquisition, the outputs requested, and
information kept between solves.

For structured systems, a sparse full-rank
base with a few controlled corrections allows scalar estimation and
sampling without forming a dense inverse or normalizing small correction
columns, so cone dimension alone does not determine query complexity.
For explicit directions, eccentricity profiles and latent sparsity give
different classical solve bounds; their assumptions must hold at the
iterates used, and their residual guarantees must match the
outer method. The parameter regimes in which our upper and lower bounds
do not match remain open.
